% ============================================================================
% APPENDICES — INDEX_INSIGHT
% ============================================================================
% Insert this block at the end of the thesis, after the Conclusion.
% Required packages in the preamble (if they are not already loaded):
% \usepackage{booktabs}
% \usepackage{longtable}
% \usepackage{array}
% ============================================================================

\appendix

\chapter{Corpus Sources and Metadata Schema}

\section{Composition of the final corpus}

Table~\ref{tab:appendix-corpus-sources} provides a compact description of the
six source collections retained in the final corpus. The corpus combines
documents supplied by the SIAGE Chair with materials collected from institutional,
research, associative, and territorial websites. Each retained document is kept
as an individual text record and remains linked to its source collection.

\begin{table}[H]
\centering
\begin{tabular}{lrp{8cm}}
\toprule
Source & Records & Main type of material \\
\midrule
CNAV & 42 & Institutional documents and prevention-related initiatives. \\
Fondation de France & 43 & Descriptions of associative and territorial initiatives. \\
SIAGE action sheets & 60 & Internal action sheets supplied by the host institution. \\
HCFEA & 23 & Institutional reports and public-action resources on ageing. \\
IReSP Autonomie & 131 & Research-project descriptions in the field of autonomy. \\
VADA & 854 & Descriptions of local age-friendly initiatives. \\
\midrule
Total & 1,153 & Heterogeneous French-language documentary corpus. \\
\bottomrule
\end{tabular}
\caption{Final corpus used for the Index\_Insight pipeline.}
\label{tab:appendix-corpus-sources}
\end{table}

The collection is designed for documentary exploration, not for statistical
representation of all French initiatives related to ageing. In particular, VADA
accounts for 854 records (approximately 74\%) and therefore influences the
observed distributions of themes, initiative types, and actors.

\section{Metadata schema}

Table~\ref{tab:appendix-schema} documents the JSON schema supplied to the
local language model. Empty strings or empty lists represent information that is
not supported by the source text. The schema separates factual metadata from
short textual summaries in order to preserve both filtering possibilities and
documentary context.

\begin{longtable}{p{3.3cm}p{2.5cm}p{8cm}}
\caption{Metadata fields used in the LLM-assisted extraction.}
\label{tab:appendix-schema}\\
\toprule
Field & Format & Purpose \\
\midrule
\endfirsthead
\toprule
Field & Format & Purpose \\
\midrule
\endhead
\bottomrule
\endfoot
\texttt{titre} & String & Title of the document or initiative. \\
\texttt{territoire} & String & Place explicitly associated with the initiative. \\
\texttt{echelle} & List & Territorial scale of the initiative. \\
\texttt{public\_vise} & Controlled string & Main target audience. \\
\texttt{nature\_initiative} & Short string & Main type of initiative. \\
\texttt{description\_generale} & String & Concise description of the initiative. \\
\texttt{contexte} & String & Contextual elements motivating the initiative. \\
\texttt{problematique} & String & Problem, need, or issue addressed. \\
\texttt{solution\_envisagee} & String & Proposed action or response. \\
\texttt{objectifs} & String & Stated objectives. \\
\texttt{porteur\_initiative} & String & Organisation responsible for, organising, or piloting the initiative. \\
\texttt{date} & List & Explicit or clearly supported temporal expressions. \\
\texttt{thematique} & List & Short thematic labels. \\
\texttt{source\_site} & String & Name of the source collection, retained for provenance. \\
\end{longtable}

\section{Controlled vocabularies}

The prompt constrains two fields to fixed vocabularies, thereby limiting
spelling variation and facilitating evaluation. The target-audience labels are
French because the source corpus and the extraction prompt are in French.

\begin{table}[H]
\centering
\begin{tabular}{p{3.6cm}p{10cm}}
\toprule
Field & Permitted values or constraints \\
\midrule
\texttt{public\_vise} & \textit{aînés}; \textit{aînés et autres}; \textit{professionnels}; \textit{autre}; \textit{public non connu}. \\
\texttt{echelle} & One or more of: \textit{quartier}; \textit{communale}; \textit{intercommunale}; \textit{départementale}; \textit{régionale}; \textit{nationale}; \textit{internationale}. \\
\texttt{date} & A list of explicit temporal expressions, preserved in their source form when possible (for example, \textit{2019}, \textit{depuis 2019}, or \textit{mars 2016}). \\
\texttt{thematique} & A list of short, homogeneous labels. The prompt proposes examples such as \textit{santé}, \textit{mobilité}, \textit{lien social}, \textit{numérique}, \textit{habitat}, \textit{prévention}, \textit{intergénérationnel}, \textit{isolement}, and \textit{participation citoyenne}. \\
\bottomrule
\end{tabular}
\caption{Controlled vocabularies and normalisation constraints.}
\label{tab:appendix-vocabularies}
\end{table}

\chapter{LLM Extraction Configuration and Prompt Specification}

\section{Inference configuration}

The VADA subset was processed locally through the Ollama chat API with the
\texttt{qwen2.5:1.5b} model. Table~\ref{tab:appendix-llm-configuration}
records the configuration used by the extraction script. Local inference was
chosen to maintain control over the data-processing workflow and to avoid
sending source texts to an external commercial API.

\begin{table}[H]
\centering
\begin{tabular}{ll}
\toprule
Parameter & Value \\
\midrule
Model & \texttt{qwen2.5:1.5b} \\
Inference service & Ollama local chat API \\
Temperature & 0 \\
Context window & 8,192 tokens \\
Maximum generated tokens & 700 \\
Input truncation threshold & 2,500 characters \\
Output format & Strict JSON \\
\bottomrule
\end{tabular}
\caption{Configuration of the local LLM extraction used for the VADA subset.}
\label{tab:appendix-llm-configuration}
\end{table}

The zero temperature setting aims to make identical executions more stable.
The 2,500-character input limit is an operational compromise and a recognised
limitation: information occurring later in a long document may not be available
to the model.

\section{Prompt specification}

The prompt was written in French to match the source language. Its central
instructions are reproduced below. The complete executable version is preserved
in the project script \texttt{Scripts/step9\_llm\_extract\_database.py}.

\begin{quote}\small
\textbf{Output constraint.} Return only strictly valid JSON, with no Markdown,
commentary, or text outside the JSON object. Use exactly the following keys:
\texttt{titre}, \texttt{territoire}, \texttt{echelle}, \texttt{public\_vise},
\texttt{nature\_initiative}, \texttt{description\_generale},
\texttt{contexte}, \texttt{problematique}, \texttt{solution\_envisagee},
\texttt{objectifs}, \texttt{porteur\_initiative}, \texttt{date},
\texttt{thematique}, and \texttt{source\_site}.

\medskip
\textbf{General rule.} Do not invent information. An inference is permitted
only when it is strongly implicit in the source text. When a value is absent or
cannot be inferred clearly, return an empty string or an empty list. Keep the
source identifier unchanged.

\medskip
\textbf{Field constraints.} For \texttt{public\_vise}, return exactly one
label from the controlled vocabulary in Table~\ref{tab:appendix-vocabularies}.
For \texttt{echelle}, return a list containing only the permitted territorial
scales. For \texttt{date}, identify all explicit temporal expressions and do
not force approximate expressions into ISO format. For
\texttt{porteur\_initiative}, identify the actor that carries, organises,
launches, pilots, or implements the initiative; leave the field empty when no
such actor is identifiable. For \texttt{thematique}, use short and homogeneous
labels and avoid duplicates caused by capitalisation.
\end{quote}

The model is additionally instructed to identify the main initiative type when
several descriptions are possible, and to formulate the field
\texttt{problematique} as a need or issue rather than merely repeating an
objective. Raw model answers are archived when available, which makes it
possible to compare an unexpected record with both the prompt and the original
response.

\chapter{Reference Annotation Protocol}

\section{Sampling design}

The reference evaluation uses a reproducible stratified random sample of 60
documents. Ten documents were drawn from each of the six source collections.
This allocation was selected to expose possible source-specific weaknesses,
including in the smaller collections; it is therefore not proportional to the
composition of the full corpus.

\begin{table}[H]
\centering
\begin{tabular}{lr}
\toprule
Source collection & Documents in reference sample \\
\midrule
CNAV & 10 \\
Fondation de France & 10 \\
SIAGE action sheets & 10 \\
HCFEA & 10 \\
IReSP Autonomie & 10 \\
VADA & 10 \\
\midrule
Total & 60 \\
\bottomrule
\end{tabular}
\caption{Stratified composition of the human reference sample.}
\label{tab:appendix-evaluation-sample}
\end{table}

For every sampled document, the reviewer compared the LLM output with the
cleaned source text. Seven fields were examined: \texttt{territoire},
\texttt{echelle}, \texttt{public\_vise}, \texttt{nature\_initiative},
\texttt{porteur\_initiative}, \texttt{date}, and \texttt{thematique}.

\section{Reference values and judgement labels}

For the controlled fields \texttt{public\_vise}, \texttt{echelle}, and
\texttt{nature\_initiative}, the reviewer records a gold value even if the
model prediction is wrong. For the other fields, the reviewer records a short
reference value when it is possible to establish one from the source. A field
that is absent from the source is recorded as \texttt{ABSENT}.

\begin{longtable}{p{3cm}p{11cm}}
\caption{Definitions of the human judgement labels.}
\label{tab:appendix-judgement-labels}\\
\toprule
Label & Definition \\
\midrule
\endfirsthead
\toprule
Label & Definition \\
\midrule
\endhead
\bottomrule
\endfoot
\textit{correct} & The predicted value is exact and explicitly supported by the source text. \\
\textit{partiel} & The value is broadly supported but incomplete, overly vague, or contains a debatable element. \\
\textit{incorrect} & The value is wrong, unsupported, or invented. \\
\textit{absence correcte} & The model returned no value and the source does not provide enough information to fill the field reliably. This is treated as an accepted prediction. \\
\textit{non évaluable} & The source is too ambiguous, incomplete, or degraded for the reviewer to establish a reliable reference value. This case is reported separately and excluded from the accepted-rate denominator. \\
\end{longtable}

When a value is only strongly implicit, it is accepted as correct only if an
independent reader could reasonably reach the same conclusion from the source
text. In doubtful cases, the protocol favours the \textit{partiel} label and
requires a short explanatory note.

\section{Reported measures}

The categorical fields are reported with micro precision, recall, micro F1,
macro F1, and exact-match rate. For the four remaining fields, the evaluation
reports the number of evaluable cases; the counts of accepted, partial,
incorrect, and non-evaluable judgments; and the accepted rate calculated over
evaluable documents only. Accepted predictions combine \textit{correct} and
\textit{absence correcte} labels.

\chapter{Reproducibility Materials}

The project is organised so that collection, cleaning, extraction, evaluation,
and analysis remain separable. Table~\ref{tab:appendix-materials} identifies
the main files and directories required to reproduce or audit the work. Raw
source documents and downloaded PDFs are not reproduced in this thesis because
they originate from third-party websites; their provenance is retained in the
corpus indexes and in the structured records.

\begin{longtable}{p{4.5cm}p{9.5cm}}
\caption{Main reproducibility materials in the project repository.}
\label{tab:appendix-materials}\\
\toprule
Material & Role \\
\midrule
\endfirsthead
\toprule
Material & Role \\
\midrule
\endhead
\bottomrule
\endfoot
\texttt{Scripts/} & Source-specific collection scripts, preprocessing modules, relevance scoring, LLM extraction, evaluation sampling, scoring, and graph generation. \\
\texttt{Database/} & Raw text records collected or supplied for the corpus. \\
\texttt{Cleaned\_Database/} & Normalised text records used as input to subsequent processing. \\
\texttt{LLM\_Tagged\_Database/} & JSON records generated by the local LLM, together with archived raw answers where available. \\
\texttt{Evaluation/llm\_annotation\_template.xlsx} & Blank workbook used for the manual reference annotation. \\
\texttt{Evaluation/annotation\_guide.md} & Detailed operational annotation instructions. \\
\texttt{Scripts/create\_llm\_evaluation\_sample.py} & Reproducible stratified sampling procedure. \\
\texttt{Scripts/score\_llm\_annotations.py} & Computation of the reported evaluation measures. \\
\texttt{Evaluation/llm\_evaluation\_results.csv} & Machine-readable table from which the evaluation results were produced. \\
\texttt{outputs\_graphs/} & Tables and visual outputs used for descriptive exploration. \\
\end{longtable}

The version of each script, the model identifier, the prompt specification, and
the inference configuration should be archived with the corresponding corpus
release. This information is necessary to distinguish changes caused by source
updates, preprocessing changes, or model and prompt revisions.
